% ============================================================================
% scMetaIntel-Hub: Article Plan & Benchmarking Design
% Version 0.3.0 — 2026-03-26
% ============================================================================
\documentclass[11pt,a4paper]{article}

% --- Packages ---
\usepackage[utf8]{inputenc}
\usepackage[T1]{fontenc}
\usepackage{geometry}
\geometry{margin=1in}
\usepackage{booktabs}
\usepackage{enumitem}
\usepackage{hyperref}
\usepackage{xcolor}
\usepackage{graphicx}
\usepackage{amsmath}
\usepackage{multirow}
\usepackage{longtable}
\usepackage{caption}

\hypersetup{colorlinks=true, linkcolor=blue!70!black, citecolor=green!50!black, urlcolor=blue!60!black}

\title{%
  \textbf{scMetaIntel-Hub}: Benchmarking Local LLMs for \\
  Single-Cell Metadata Intelligence \\[6pt]
  \large Article Plan \& Benchmarking Design (v0.3.0)}
\author{Zeyu Fu}
\date{March 2026}

\begin{document}
\maketitle
\tableofcontents
\newpage

% ============================================================================
\section{Article Positioning}
% ============================================================================

\subsection{Central Claim}

We present \textbf{scMetaIntel-Hub}, a fully local, privacy-preserving
pipeline that converts natural language queries into grounded, cited answers
over single-cell genomics metadata --- without requiring gene expression
matrices or cloud API access. Through systematic benchmarking of 21 LLM
configurations across 5 task dimensions, we show that mid-size local models
(9B parameters) match or exceed larger models on metadata-centric tasks when
paired with the right retrieval strategy.

\subsection{Key Contributions}

\begin{enumerate}[leftmargin=2em]
  \item \textbf{Metadata-only intelligence}: A complete pipeline (enrich
        $\to$ ontology $\to$ embed $\to$ retrieve $\to$ answer) that operates
        exclusively on GEO metadata and PubMed abstracts --- no expression
        matrices, no cloud dependencies.
  \item \textbf{Comprehensive local LLM benchmark}: 21 model configurations
        (5 families: Qwen3/3.5, LLaMA~3.1, Gemma~3, Mistral, Command-R)
        evaluated across 5 tasks (parsing, extraction, ontology normalization,
        answer generation, speed) on a curated 2{,}189-study corpus with 90
        evaluation queries.
  \item \textbf{Optimal pipeline recipe}: Bench-validated configuration
        achieving 100\% citation recall, 100\% grounding, 93.1\% query parse
        exact match, and 0.428 hybrid retrieval R@50 --- all running locally
        on a single laptop GPU (RTX~5090, 25\,GB VRAM).
  \item \textbf{Curated evaluation resource}: 2{,}189 enriched GSE metadata
        JSONs spanning 43 organisms, 90 evaluation queries across 7
        categories, and a reproducible cleaning pipeline for ground truth
        quality.
\end{enumerate}

\subsection{Target Venue}

\begin{itemize}
  \item Primary: \textit{Bioinformatics} (Applications Note or Original Paper)
  \item Alternative: \textit{Nucleic Acids Research} (Web Server / Database issue),
        \textit{Briefings in Bioinformatics}, or \textit{Genome Biology} (Tools)
\end{itemize}

% ============================================================================
\section{Article Structure}
% ============================================================================

% ---------------------------------------------------------------------------
\subsection{Title Options}

\begin{enumerate}
  \item \textit{scMetaIntel-Hub: Benchmarking Local LLMs for Privacy-Preserving
        Single-Cell Dataset Discovery}
  \item \textit{Local LLMs for Single-Cell Metadata Intelligence: A Systematic
        Benchmark of 21 Models Across 5 Tasks}
  \item \textit{From Query to Citation: A Fully Local Pipeline for Grounded
        Single-Cell Dataset Search}
\end{enumerate}

% ---------------------------------------------------------------------------
\subsection{Abstract (Target: 250 words)}

\paragraph{Structure:}
\textbf{Background} (metadata explosion, need for intelligent search,
privacy concerns with cloud APIs) $\to$ \textbf{Method} (local LLM pipeline,
hybrid retrieval, 5-task benchmark on 2{,}189 GEO studies) $\to$
\textbf{Results} (top model picks, 100\% citation recall, bigger $\neq$
better finding) $\to$ \textbf{Conclusion} (practical recipe, open resource).

% ---------------------------------------------------------------------------
\subsection{Section 1: Introduction (1.5 pages)}

\begin{enumerate}[leftmargin=2em]
  \item Single-cell genomics growth: 2{,}000+ new GEO datasets per year;
        metadata is the primary discovery interface.
  \item Problem: existing search (GEO text search, CellxGene) is keyword-only;
        no semantic understanding, no structured constraint parsing, no
        cited answers.
  \item LLM opportunity: local models can parse, extract, normalize, and
        generate --- but which model, for which task?
  \item Gap: no systematic benchmark of local LLMs on single-cell metadata
        tasks. Existing LLM-for-bio benchmarks (e.g., BioASQ, PubMedQA)
        focus on literature QA, not structured metadata retrieval.
  \item Contribution summary (4 points from \S1.2).
\end{enumerate}

% ---------------------------------------------------------------------------
\subsection{Section 2: Methods (3--4 pages)}

\subsubsection{2.1 Data Collection and Curation}

\begin{itemize}
  \item \textbf{Corpus construction}: 2{,}189 GEO Series (GSE) metadata
        documents enriched from GEO SOFT + PubMed E-utilities.
  \item \textbf{Schema}: \texttt{gse\_id}, \texttt{title}, \texttt{summary},
        \texttt{organism}, \texttt{tissues}, \texttt{diseases},
        \texttt{cell\_types}, \texttt{document\_text}, \texttt{pubmed},
        \texttt{samples}.
  \item \textbf{Cleaning pipeline}: (1) remove non-diseases (healthy,
        control, etc.), (2) migrate disease terms from tissue field, (3)
        extract diseases/cell-types from text when annotations are empty,
        (4) fix encoding artifacts (\texttt{scripts/clean\_ground\_truth.py}).
  \item \textbf{Coverage}: 83.8\% tissues, 44.8\% diseases, 72.6\%
        cell types, 69.0\% PubMed, 100\% document text. 43 organisms.
  \item \textbf{Table 1}: Corpus statistics by organism, disease category,
        modality.
\end{itemize}

\subsubsection{2.2 Pipeline Architecture}

\begin{itemize}
  \item \textbf{Figure 1}: Architecture diagram. Query $\to$
        Parse (LLM) $\to$ Embed (HF model) $\to$
        Retrieve (Qdrant hybrid dense+sparse) $\to$
        Context Assembly ($k$=3, structured format) $\to$
        Answer Generation (LLM with citations).
  \item Ontology normalization: 3-layer (exact/synonym $\to$ BioLORD
        embedding $\to$ LLM fallback) against CL, UBERON, MONDO.
  \item All components are local: Ollama (LLMs), HuggingFace
        (embeddings/rerankers), Qdrant (vector store).
\end{itemize}

\subsubsection{2.3 Evaluation Queries}

\begin{itemize}
  \item 90 queries across 7 categories:
        basic (19), disease (23), cell\_type (12),
        multi\_constraint (15), natural\_language (8),
        modality (7), organism (6).
  \item 3 difficulty levels: easy (22), medium (37), hard (21).
  \item 227 unique expected GSEs (10.4\% corpus coverage).
  \item \textbf{Table 2}: Query category distribution with examples.
\end{itemize}

\subsubsection{2.4 Benchmark Design (5 Levels)}

\paragraph{Level 2 --- Embedding comparison (16 models)}

\begin{itemize}
  \item Models: 10 general (mxbai, bge-m3, nomic, gte, stella, qwen3-embed
        $\times$3, jina, snowflake) + 6 biomedical (biolord,
        biomedbert $\times$2, medcpt $\times$2, specter2).
  \item 4 sub-tasks: semantic clustering (silhouette), retrieval accuracy
        (P@5/10, R@50, MRR, nDCG@10), ontology matching (recall@1/5),
        encoding speed (tok/s).
\end{itemize}

\paragraph{Level 3 --- Retrieval strategy comparison (6 strategies)}

\begin{itemize}
  \item Strategies: dense, sparse, hybrid, hybrid+filter, hybrid+rerank,
        hybrid+filter+rerank.
  \item Best embedding from Level~2 (mxbai-embed-large).
  \item Metrics: P@5, P@10, R@50, MRR, nDCG@10, latency.
\end{itemize}

\paragraph{Level 4 --- LLM comparison (21 configurations)}

\begin{itemize}
  \item Models: 5 families, 7B--32B, Q4\_K\_M and Q8\_0 quantizations,
        thinking-mode ablation where supported.
  \item 5 sub-tasks:
    \begin{enumerate}
      \item[A.] \textbf{Query Parsing}: NL $\to$ structured JSON
            (organism/tissue/disease/assay/cell\_type). Metric: exact match, field accuracy.
      \item[B.] \textbf{Metadata Extraction}: title+summary $\to$
            tissues/diseases/cell\_types arrays. Metric: per-field P/R/F1.
      \item[C.] \textbf{Ontology Normalization}: raw terms $\to$
            UBERON/CL/MONDO IDs. Metric: accuracy, F1.
      \item[D.] \textbf{Answer Generation}: query + context $\to$ cited
            answer. Metric: citation precision/recall, grounding rate.
      \item[E.] \textbf{Inference Speed}: tok/s throughput.
    \end{enumerate}
  \item Composite score: weighted combination of all 5 tasks.
  \item \textbf{Table 3}: Full LLM comparison matrix (21 rows $\times$
        5 tasks + composite).
\end{itemize}

\paragraph{Level 5 --- Context optimization}

\begin{itemize}
  \item Variables: $k$ (1--10), format (full/structured/minimal), system
        prompt variants.
  \item Advanced strategies: MMR diversity, recency reorder, ontology
        expansion, multi-step retrieval.
\end{itemize}

\paragraph{Level 7 --- End-to-end pipeline}

\begin{itemize}
  \item Full pipeline: query $\to$ parse $\to$ retrieve $\to$ answer.
  \item Metrics: end-to-end latency, citation quality, answer correctness.
\end{itemize}

\subsubsection{2.5 Hardware and Environment}

\begin{itemize}
  \item GPU: NVIDIA RTX 5090 Laptop (25\,GB VRAM), CUDA 13.0.
  \item LLM serving: Ollama with \texttt{num\_ctx=4096}.
  \item Embeddings: HuggingFace + PyTorch 2.9.0 (conda \texttt{dl} env).
  \item Vector store: Qdrant (local persistent).
  \item OS: Linux 6.14.
\end{itemize}

% ---------------------------------------------------------------------------
\subsection{Section 3: Results (3--4 pages)}

\subsubsection{3.1 Embedding Results}

\begin{itemize}
  \item \textbf{Finding}: mxbai-embed-large dominates with R@50=0.410 dense,
        0.428 hybrid+filter.
  \item \textbf{Finding}: Biomedical embedders (biolord, medcpt) are strong
        at ontology matching but weaker at general retrieval.
  \item \textbf{Figure 2}: Radar chart --- general vs biomedical embedders
        across 4 sub-tasks.
\end{itemize}

\subsubsection{3.2 Retrieval Strategy Results}

\begin{itemize}
  \item \textbf{Finding}: hybrid+filter is optimal (R@50=0.428, MRR=0.190,
        82\,ms). Reranking \textit{hurts} (-MRR, +110\,ms) because
        the non-biomedical reranker (ms-marco) degrades domain relevance.
  \item \textbf{Table 4}: 6 strategies $\times$ metrics. Breakdown by query
        difficulty.
\end{itemize}

\subsubsection{3.3 LLM Benchmark Results}

\begin{itemize}
  \item \textbf{Key Finding 1}: \textit{Bigger is NOT better} --- qwen3.5-9b
        (0.800 composite) beats qwen3-32b and command-r-35b.
  \item \textbf{Key Finding 2}: Qwen3.5 family dominates top 3 (27b, 9b,
        9b-q8). Thinking mode helps parsing (+3\%) but not answer
        generation.
  \item \textbf{Key Finding 3}: Best model per task diverges ---
        llama3.1-8b for parsing (93.1\% EM), qwen2.5-7b for extraction
        (0.380 F1), qwen3.5-9b for answers (100\% recall).
  \item \textbf{Key Finding 4}: 100\% answer citation recall requires
        $\geq$9B parameters.
  \item \textbf{Figure 3}: Composite score vs model size scatter plot
        (annotated with model names, color by family).
  \item \textbf{Figure 4}: Per-task heatmap (21 models $\times$ 5 tasks).
  \item \textbf{Table 5}: Top-5 model comparison with per-task breakdown.
\end{itemize}

\subsubsection{3.4 Context Optimization Results}

\begin{itemize}
  \item $k$=3 structured format achieves 100\% precision, 91.1\% recall,
        168 tokens average context.
  \item Diminishing returns beyond $k$=5.
  \item \textbf{Figure 5}: Precision/recall curve vs $k$.
\end{itemize}

\subsubsection{3.5 End-to-End Pipeline Performance}

\begin{itemize}
  \item Full pipeline latency breakdown: parse (X\,ms) + retrieve (82\,ms)
        + answer (Y\,ms).
  \item \textbf{Table 6}: Optimal pipeline recipe card.
\end{itemize}

% ---------------------------------------------------------------------------
\subsection{Section 4: Discussion (1.5 pages)}

\begin{enumerate}[leftmargin=2em]
  \item \textbf{Task-specific routing matters}: No single model wins all
        tasks. A router that dispatches llama3.1-8b for parsing and
        qwen3.5-9b for answers outperforms any monolithic choice.
  \item \textbf{Quantization trade-offs}: Q8\_0 vs Q4\_K\_M yields
        negligible quality difference at 1.4$\times$ speed cost.
        Q3\_K\_M only viable for standard-RoPE models without CPU spill.
  \item \textbf{Reranking paradox}: Generic rerankers degrade
        domain-specific retrieval. Bio-domain rerankers are needed.
  \item \textbf{Metadata quality is the real bottleneck}: Disease annotations
        are sparse (44.8\% after aggressive extraction). We release the
        cleaning pipeline to help the community.
  \item \textbf{Limitations}: (a) GEO-centric (no CellxGene, HCA, or other
        portals), (b) English-only, (c) metadata only --- no expression
        data integration, (d) single-GPU evaluation.
  \item \textbf{Comparison with related work}: contrast with scBERT,
        GeneFormer, CellPLM (expression-based), PubMedQA, BioASQ
        (literature QA), and CellxGene/GEO metadata search
        (keyword-only).
\end{enumerate}

% ---------------------------------------------------------------------------
\subsection{Section 5: Conclusion (0.5 page)}

\begin{enumerate}[leftmargin=2em]
  \item Summarize optimal recipe: llama3.1-8b (parse) + mxbai-embed-large
        (embed) + hybrid+filter (retrieve) + qwen3.5-9b (answer).
  \item \textbf{Practical message}: A 9B model on a laptop GPU achieves
        state-of-the-art metadata intelligence without cloud APIs.
  \item \textbf{Resource release}: Corpus (2{,}189 GSEs), 90 eval queries,
        benchmark scripts, cleaning pipelines --- all open-source.
\end{enumerate}

% ============================================================================
\section{Figures and Tables Plan (8 total)}
% ============================================================================

\begin{longtable}{@{}clp{7cm}@{}}
\toprule
\# & Type & Description \\
\midrule
1 & Figure & Pipeline architecture diagram \\
2 & Figure & Embedding radar chart (general vs biomedical) \\
3 & Figure & Composite score vs model size scatter plot \\
4 & Figure & Per-task heatmap (21 models $\times$ 5 tasks) \\
5 & Figure & Precision/recall vs context $k$ \\
6 & Table  & Corpus statistics (organism, disease category, modality) \\
7 & Table  & Retrieval strategy comparison (6 strategies) \\
8 & Table  & LLM comparison matrix (top-5 with per-task breakdown) \\
\bottomrule
\end{longtable}

% ============================================================================
\section{Supplementary Materials}
% ============================================================================

\begin{itemize}
  \item \textbf{Supp Table S1}: Full 21-model LLM comparison (all metrics).
  \item \textbf{Supp Table S2}: Full 16-model embedding comparison.
  \item \textbf{Supp Table S3}: 90 evaluation queries with expected GSEs.
  \item \textbf{Supp Table S4}: Corpus disease/tissue/cell-type distributions.
  \item \textbf{Supp Table S5}: Quantization ablation (Q4\_K\_M vs Q8\_0 vs Q3\_K\_M).
  \item \textbf{Supp Figure S1}: Thinking-mode ablation (think=True vs False).
  \item \textbf{Supp Figure S2}: Context management strategies comparison.
  \item \textbf{Supp Figure S3}: Retrieval performance by query difficulty.
\end{itemize}

% ============================================================================
\section{Re-Benchmarking Plan (Pre-Article)}
% ============================================================================

The corpus has changed (1{,}994 $\to$ 2{,}189 studies; 67 $\to$ 90 eval
queries; cleaned ground truth). All previous benchmark results were computed
on the old corpus. The following re-benchmarking sequence is required before
writing the article.

\subsection{Phase 1: Infrastructure (no GPU needed)}

\begin{enumerate}[leftmargin=2em]
  \item \textbf{Re-index Qdrant}: Re-embed all 2{,}189 documents with
        mxbai-embed-large into a fresh Qdrant collection.
        \begin{itemize}
          \item Script: \texttt{python -m scmetaintel embed}
          \item Duration: $\sim$5--10 min
          \item Validates: new documents are indexed; old index is replaced.
        \end{itemize}
  \item \textbf{Verify eval query GSEs}: Confirm all 227 expected GSEs
        from the 90 queries exist in the ground truth and in Qdrant.
\end{enumerate}

\subsection{Phase 2: Embedding Re-Bench (bench 02)}

\begin{enumerate}[leftmargin=2em]
  \item Run \texttt{python benchmarks/02\_bench\_embeddings.py} on the
        2{,}189-document corpus with 90 queries.
  \item Expected: mxbai-embed-large likely still wins, but numbers will
        shift (more documents $\to$ retrieval is harder).
  \item Duration: $\sim$2--3 hours (16 models).
\end{enumerate}

\subsection{Phase 3: Retrieval Re-Bench (bench 03)}

\begin{enumerate}[leftmargin=2em]
  \item Run \texttt{python benchmarks/03\_bench\_retrieval.py} on the fresh
        Qdrant index with 90 queries.
  \item Compare 6 strategies on the expanded query set.
  \item Special attention: new disease/modality queries may shift
        hybrid+filter advantage.
  \item Duration: $\sim$30 min.
\end{enumerate}

\subsection{Phase 4: LLM Re-Bench (bench 04)}

\begin{enumerate}[leftmargin=2em]
  \item Run \texttt{python benchmarks/04\_bench\_llm.py} on the cleaned
        ground truth + 90 queries.
  \item 21 configurations $\times$ 5 tasks $\times$ thinking ablation.
  \item Critical: uses cleaned gold truth (diseases went from $\sim$9\%
        $\to$ 45\% annotated), so extraction scores \textit{will} change.
  \item Duration: $\sim$6--10 hours (GPU-bound, one model at a time).
\end{enumerate}

\subsection{Phase 5: Context \& E2E Re-Bench (bench 05/07)}

\begin{enumerate}[leftmargin=2em]
  \item Run \texttt{python benchmarks/05\_bench\_context.py} and
        \texttt{07\_bench\_e2e.py}.
  \item Validates optimal $k$ and end-to-end pipeline on expanded queries.
  \item Duration: $\sim$1--2 hours.
\end{enumerate}

\subsection{Phase 6: Results Analysis \& Figure Generation}

\begin{enumerate}[leftmargin=2em]
  \item Load all re-bench JSON results.
  \item Generate all figures (Figs~1--5) and tables (Tables~1--8).
  \item Compute final composite scores and model rankings.
  \item Write results prose based on actual re-bench numbers.
\end{enumerate}

\subsection{Re-Bench Schedule Estimate}

\begin{tabular}{@{}lll@{}}
\toprule
Phase & Duration & Dependency \\
\midrule
1. Re-index Qdrant         & 10 min    & None \\
2. Embedding bench (02)    & 2--3 hr   & Phase 1 \\
3. Retrieval bench (03)    & 30 min    & Phase 1 \\
4. LLM bench (04)          & 6--10 hr  & None (independent) \\
5. Context \& E2E (05/07)  & 1--2 hr   & Phases 1, 3 \\
6. Analysis \& figures     & 2--3 hr   & Phases 2--5 \\
\midrule
\textbf{Total}             & \textbf{$\sim$12--19 hr}  & Phases 2+3 parallelizable \\
\bottomrule
\end{tabular}

\bigskip
\noindent\textit{Phases 2+3 can run in parallel with Phase 4. Total wall-clock
$\approx$ 10--14 hours if parallelized.}

% ============================================================================
\section{Risks and Mitigations}
% ============================================================================

\begin{longtable}{@{}p{5cm}p{3cm}p{5cm}@{}}
\toprule
Risk & Likelihood & Mitigation \\
\midrule
Re-bench results differ significantly from old results &
  High (expected) &
  Article uses only re-bench numbers; old results archived \\
Retrieval R@50 drops with larger corpus &
  Medium &
  Expected (more distractors); hybrid+filter advantage should increase \\
Extraction F1 changes due to cleaned gold &
  High &
  Cleaned gold is correct; old gold was polluted \\
Some eval query GSEs missing from Qdrant after re-index &
  Low &
  Verification step in Phase~1 catches this \\
GPU memory issues with 21 models &
  Low &
  \texttt{--include-spill} flag + \texttt{num\_ctx=4096} already proven \\
\bottomrule
\end{longtable}

% ============================================================================
\section{Current Project State Summary}
% ============================================================================

\begin{tabular}{@{}lr@{}}
\toprule
Metric & Value \\
\midrule
Ground truth GSE files       & 2{,}189 \\
Eval queries                 & 90 \\
Unique expected GSEs         & 227 \\
Query categories             & 7 \\
Organisms represented        & 43 \\
Disease coverage (cleaned)   & 44.8\% \\
Cell type coverage (cleaned) & 72.6\% \\
Tissue coverage (cleaned)    & 83.8\% \\
LLM configurations benched   & 21 \\
Embedding models benched     & 16 \\
Retrieval strategies benched & 6 \\
\bottomrule
\end{tabular}

\end{document}
